\chapter{Les nombres réels et leurs opérations}
\label{workbook:chapter:01}
\section{Les grandes familles de nombres}
\label{workbook:section:01.01}
\begin{exerciseblock}{Les grandes familles de nombres}
\item Classer chacun des nombres suivants dans le plus petit ensemble parmi $\N$, $\Z$, $\Q$ et $\R$ : $-7$, $\frac{18}{6}$, $0{,}125$, $\sqrt{49}$, $\sqrt{3}$, $\pi-3$.
\item Donner un exemple montrant que la somme de deux nombres irrationnels peut être rationnelle, puis un exemple où elle reste irrationnelle.
\item Représenter sur une droite réelle les ensembles $A=[-2,3[$ et $B=]1,5]$, puis déterminer $A\cap B$, $A\cup B$ et $A\setminus B$.
\end{exerciseblock}
\section{Priorité des opérations et signes}
\label{workbook:section:01.02}
\begin{exerciseblock}{Priorité des opérations et signes}
\item Calculer en détaillant l'ordre des opérations : $-3^2+4(2-5)^2\div 6$.
\item Comparer $(-2)^4$, $-2^4$, $(-2)^{-2}$ et $-2^{-2}$. Expliquer le rôle des parenthèses.
\item Une formule donne $E=-ab^2/c$. Déterminer le signe de $E$ lorsque $a<0$, $b\neq0$ et $c<0$, sans choisir de valeurs numériques.
\end{exerciseblock}
\section{Fractions et nombres décimaux}
\label{workbook:section:01.03}
\begin{exerciseblock}{Fractions et nombres décimaux}
\item Effectuer et simplifier : $\frac{5}{12}-\frac{7}{18}+\frac{1}{9}$.
\item Écrire sous forme de fraction irréductible : $0{,}375$, $1{,}\overline{27}$ et $0{,}1\overline{6}$.
\item Une recette utilise $\frac34$ de tasse pour $5$ portions. Quelle quantité faut-il pour $8$ portions? Donner une valeur exacte puis décimale.
\end{exerciseblock}
\section{Puissances, notation scientifique et racines}
\label{workbook:section:01.04}
\begin{exerciseblock}{Puissances, notation scientifique et racines}
\item Simplifier sans calculatrice : $\dfrac{(2^3)^4\,2^{-5}}{2^2}$ et $\dfrac{3^5\,9^{-1}}{27}$.
\item Écrire en notation scientifique puis calculer : $(6{,}0\times10^{-4})(2{,}5\times10^7)$.
\item Simplifier $\sqrt{180}-2\sqrt{20}+\sqrt{45}$, puis expliquer pourquoi $\sqrt{x^2}=|x|$ et non toujours $x$.
\end{exerciseblock}
\section{La droite réelle comme modèle unificateur}
\label{workbook:section:01.05}
\begin{exerciseblock}{La droite réelle comme modèle unificateur}
\item Résoudre graphiquement puis algébriquement $|x-2|\leq 3$.
\item Décrire par une inégalité l'ensemble des réels situés à moins de $0{,}4$ de $-1{,}5$.
\item Deux mesures sont $a=4{,}8$ et $b=5{,}3$. Comparer leur écart absolu et leur écart relatif par rapport à $b$.
\end{exerciseblock}
\section{Cohérence des règles de calcul}
\label{workbook:section:01.06}
\begin{exerciseblock}{Cohérence des règles de calcul}
\item À partir de $a^m/a^n=a^{m-n}$, justifier les conventions $a^0=1$ et $a^{-n}=1/a^n$ pour $a\neq0$.
\item Décider si les égalités suivantes sont toujours vraies, parfois vraies ou jamais vraies : $\sqrt{a+b}=\sqrt a+\sqrt b$; $(a+b)^2=a^2+b^2$; $\sqrt{ab}=\sqrt a\sqrt b$.
\item Construire un contre-exemple numérique à la règle erronée $\frac{a+b}{c+d}=\frac ac+\frac bd$.
\end{exerciseblock}
\section{Mesurer, estimer et contrôler un résultat}
\label{workbook:section:01.07}
\begin{exerciseblock}{Mesurer, estimer et contrôler un résultat}
\item Sans effectuer le calcul exact, encadrer $\dfrac{19{,}8\times 4{,}03}{0{,}51}$ entre deux dizaines consécutives.
\item Une vitesse est donnée par $v=d/t$ avec $d=240\ \mathrm{km}$ et $t=3{,}2\ \mathrm h$. Vérifier les unités et estimer le résultat avant de calculer.
\item Un logiciel annonce que $37\%$ de $480$ vaut $1776$. Donner deux contrôles rapides qui détectent l'erreur, puis corriger.
\end{exerciseblock}
\section*{Synthèse du chapitre}
\begin{synthesisbox}
\item Une calculatrice affiche $\sqrt{2}\approx1{,}41421356$. Expliquer pourquoi cette écriture ne transforme pas $\sqrt2$ en nombre rationnel et distinguer valeur exacte et valeur approchée.
\item Une facture de $84\,$\$ est augmentée de $15\%$, puis réduite de $15\%$. Calculer le montant final et expliquer pourquoi on ne revient pas au montant initial.
\item On mesure un rectangle $12{,}4\pm0{,}1$ cm par $8{,}1\pm0{,}1$ cm. Estimer l'aire et encadrer l'aire réelle à l'aide des valeurs extrêmes.
\end{synthesisbox}
